\chapter{Единая родительская нормировка жёсткостей
\texorpdfstring{$Q$ и $T$}{Q и T}}

\section{Возврат к условному следу}

Предыдущая глава локализовала открытый коэффициент $Z_Q$. Для его вывода
не требуется новый упругий параметр: поля $Q$ и $T$ уже живут над одним
семейным триплетом и читаются одним условным следом. Необходимо, однако,
строго сопоставить нормировки, использованные в гейте тетраэдрической
кривизны и в последующей краевой задаче.

Гейт модулярного угла спина три доказал тождество
\begin{equation}
 \frac{\tau_8(p_1Xp_1)}{\tau_8(p_1)}
 =\frac13\Tr_3X.
 \label{eq:v6-stiffness-conditioned-trace}
\end{equation}
Литературная основа условных следов и норм гильбертовых модулей собрана в
соответствующей карточке источников по угловой кривизне; общая граница
независимых упругих инвариантов $Q$ --- в карточке источников по упругой
нормировке $Q$-тензора.

\section{Кинетическая норма поля \texorpdfstring{$Q$}{Q}}

Пусть $E_a$, $a=1,\ldots,5$, --- ортонормированный по обычному следу базис
$\operatorname{Sym}_0^2(V_1)$:
\begin{equation}
 \Tr_3(E_aE_b)=\delta_{ab}.
\end{equation}
Тогда условный след~\eqref{eq:v6-stiffness-conditioned-trace} даёт
\begin{equation}
 \frac13\Tr_3(E_aE_b)=\frac13\delta_{ab}.
\end{equation}
Следовательно, в обычной фробениусовой норме
\begin{equation}
 Z_Q=\frac13.
 \label{eq:v6-canonical-ZQ}
\end{equation}

\section{Кинетическая норма поля \texorpdfstring{$T$}{T}}

Каноническая изометрия предыдущих глав имеет вид
\begin{equation}
 Z_T:V_1\longrightarrow V_2,
 \qquad
 (Z_Te_k)_{ij}=T_{ijk}.
\end{equation}
Поэтому
\begin{equation}
 \Tr_3\bigl((DZ_T)^*DZ_T\bigr)=\|DT\|^2,
\end{equation}
а тот же условный след даёт
\begin{equation}
 \frac13\Tr_3\bigl((DZ_T)^*DZ_T\bigr)
 =\frac13\|DT\|^2.
 \label{eq:v6-canonical-ZT}
\end{equation}
Машинный аудит на ортонормированных базисах получил пять и семь
собственных значений, каждое равное $1/3$ с погрешностью ниже
$2\cdot10^{-16}$. Таким образом,
\begin{equation}
 Z_Q=Z_T=\frac13,
 \qquad
 \frac{Z_Q}{Z_T}=1.
 \label{eq:v6-relative-QT-stiffness}
\end{equation}
Нового относительного веса нет.

\section{Обнаруженная несогласованность старого профиля}

Для разложения тетраэдрического вакуума уже доказано
\begin{equation}
 \|T_3\|^2=\frac{128}{81},
 \qquad
 \|T_0\|^2=\frac{160}{81}.
\end{equation}
После условной нормировки~\eqref{eq:v6-canonical-ZT} коэффициенты
пространственной кинетики обязаны быть
\begin{equation}
 A=C=\frac{128}{243},
 \qquad
 B=\frac{160}{243}.
 \label{eq:v6-corrected-vortex-kinetic-coefficients}
\end{equation}

Прежняя краевая задача использовала $128/81$ и $160/81$, то есть обычную
тензорную норму без деления на три. При этом потенциальный квадрат уже
был условно нормирован множителем $1/3$, а коэффициент связности
\begin{equation}
 G=\frac{2}{27}
\end{equation}
также был вычислен условным триплетным следом. Следовательно, прежний
профиль смешивал две нормировки и завышал только кинетику $T$ ровно в три
раза.

\begin{proposition}[Исправление без нового параметра]
Единый условный след одновременно фиксирует $Z_Q$, $Z_T$ и относительную
нормировку полей. Исправление коэффициентов
\eqref{eq:v6-corrected-vortex-kinetic-coefficients} не является новой
подгонкой: оно устраняет противоречие между уже доказанным следом и
краевой задачей.
\end{proposition}

\section{Исправленный профиль}

Краевая задача повторно решена с коэффициентами
\eqref{eq:v6-corrected-vortex-kinetic-coefficients}. Получено
\begin{align}
 b(0)&=1.0851729254\ldots,\\
 a'(0)&=2.9597500000\ldots,\\
 \mathcal T_{\mathrm{corr}}&=1.5772484656\ldots.
\end{align}
Вклады в натяжение равны
\begin{equation}
 (E_{r},E_{\theta},E_F,E_V)
 =(0.48520178,0.44610475,0.32297096,0.32297097).
\end{equation}
Относительный вириальный остаток равен $6.07\cdot10^{-9}$, а максимальный
остаток краевой задачи --- $2.0\cdot10^{-7}$.

\section{Повторная радиальная устойчивость}

Исправленный трёхпрофильный гессиан проверен на сетках от ста до пятисот
узлов. Нижнее собственное значение сходится как
\begin{equation}
 4.14140900, 4.14322268, 4.14371751,
 4.14388918, 4.14394146.
\end{equation}
На последней сетке первые восемь значений положительны:
\begin{equation}
\begin{split}
 &(4.14394146,4.68982279,4.77791339,4.92350070,\\
 &\hspace{1cm}5.12391448,5.37725068,5.68211303,6.03727402).
\end{split}
\end{equation}
Отрицательных и нулевых радиальных мод нет.

\section{Изменение статуса прежних сертификатов}

Прежние главы о профиле, радиальной, эффективной нерадиальной и полном
внутреннем гессианах сохраняются как воспроизводимая история условной
ветви. Но их численные спектры относятся к завышенной втрое кинетике
$T$ и не являются каноническими следствиями родительского следа.

\begin{nogocriterion}[Запрет переноса старого нерадиального спектра]
Нельзя переносить положительность прежнего нерадиального оператора на
исправленный профиль без повторного вычисления. Изменились профиль,
натяжение и все производные коэффициенты, хотя топологический класс
$Z_6$ сохранился.
\end{nogocriterion}

\begin{protocol}
\begin{enumerate}
 \item родительская жёсткость $Q$: \textbf{$Z_Q=1/3$};
 \item родительская жёсткость $T$: \textbf{$Z_T=1/3$};
 \item новый относительный вес: \textbf{нет};
 \item ошибка прежней кинетики $T$: \textbf{множитель три};
 \item исправленный профиль: \textbf{получен};
 \item исправленное натяжение: \textbf{$1.5772484656\ldots$};
 \item радиальная устойчивость: \textbf{проходит};
 \item исправленная нерадиальная устойчивость: \textbf{не проверена};
 \item следующий гейт: \textbf{нерадиальный гессиан исправленного вихря}.
\end{enumerate}
\end{protocol}

Машинный сертификат сохранён в
\path{s2t_v6_bosonic_defect_q_stiffness_parent_normalization_gate_results.json}.